\documentclass[a4paper,oneside,cs4size,fancyhdr]{ctexart}
\usepackage[top=1in,bottom=1in,left=1in,right=1in]{geometry} % 用于设置页面布局
\usepackage{xeCJK}  % 用于使用本地字体
\usepackage[super, square, sort&compress]{natbib}   % 处理参考文献
\usepackage{titlesec, titletoc} % 设置章节标题及页眉页脚
%\usepackage{xCJKnumb}  % 中英文数字转换
\usepackage{amssymb}
\usepackage{amsmath}    % 在公式中用\text{文本}输入中文
\usepackage{diagbox}    % 斜线表
\usepackage{multirow}   % 表格中使用多行
\usepackage{booktabs}   % 三线表
\usepackage{rotating}   % 使用sidewaystable环境旋转表格
\usepackage{tabularx}   % 定宽表格
\usepackage{graphicx}   % 处理图片
\usepackage{footnote}   % 增强的脚注功能，可添加表格脚注
\usepackage{threeparttable} % 添加真正的表格脚注,示例见README
\usepackage{hyperref}   % 添加pdf书签
\usepackage{listings}   % 添加代码高亮

% 字体设置
\setmainfont{Times New Roman}
\setsansfont[Scale=MatchLowercase,Mapping=tex-text]{DejaVu Sans}
\setmonofont[Scale=MatchLowercase]{DejaVu Sans Mono}
\setCJKmainfont[ItalicFont={Adobe Kaiti Std}, BoldFont={Adobe Heiti Std}]{Adobe Song Std}
\setCJKsansfont{Adobe Heiti Std}
\setCJKmonofont{Adobe Fangsong Std}
% \setCJKmainfont[BoldFont={FZXiaoBiaoSong-B05S}]{Adobe Song Std}
% \setCJKfamilyfont{kai}[BoldFont=Adobe Heiti Std]{Adobe Kaiti Std}
% \setCJKfamilyfont{song}[BoldFont=Adobe Heiti Std]{Adobe Song Std}
% \setCJKfamilyfont{hei}[BoldFont=Adobe Heiti Std]{Adobe Heiti Std}
% \setCJKfamilyfont{fsong}[BoldFont=Adobe Heiti Std]{Adobe Fangsong Std}
% \setCJKfamilyfont{fzwei}{FZBeiWeiKaiShu-S19S}
% \setCJKfamilyfont{fzsong}{FZXiaoBiaoSong-B05S}
% \newcommand{\kai}[1]{{\CJKfamily{kai}#1}}
% \newcommand{\hei}[1]{{\CJKfamily{hei}#1}}
% \newcommand{\fsong}[1]{{\CJKfamily{fsong}#1}}
% \setromanfont[Mapping=tex-text]{TeXGyrePagella}
% \setsansfont[Scale=MatchLowercase,Mapping=tex-text]{TeXGyrePagella}
% \setmonofont[Scale=MatchLowercase]{Courier New}
%%设置常用中文字号，方便调用
\newcommand{\erhao}{\fontsize{22pt}{\baselineskip}\selectfont}
\newcommand{\xiaoerhao}{\fontsize{18pt}{\baselineskip}\selectfont}
\newcommand{\sanhao}{\fontsize{16pt}{\baselineskip}\selectfont}
\newcommand{\xiaosanhao}{\fontsize{15pt}{\baselineskip}\selectfont}
\newcommand{\sihao}{\fontsize{14pt}{\baselineskip}\selectfont}
\newcommand{\xiaosihao}{\fontsize{12pt}{\baselineskip}\selectfont}
\newcommand{\wuhao}{\fontsize{10.5pt}{\baselineskip}\selectfont}
\newcommand{\xiaowuhao}{\fontsize{9pt}{\baselineskip}\selectfont}
\newcommand{\liuhao}{\fontsize{7.5pt}{\baselineskip}\selectfont}

% 章节标题显示方式及页眉页脚设置
% \item xCJKnumb是自己额外安装的包
% \item titleformat命令定义标题的形式
% \item titlespacing定义标题距左、上、下的距离
\titleformat{\section}{\raggedright\large\bfseries}{\thesection}{1em}{}
\titleformat{\subsection}{\raggedright\normalsize\bfseries}{\thesubsection}{1em}{}
\titlespacing{\section}{0pt}{*0}{*2}
\titlespacing{\subsection}{0pt}{*0}{*1}
% 由于默认的2em缩进不够，所以我手动调整了，但是在windows下似乎2.2就差不多了，或者是article中没有这个问题
\setlength{\parindent}{2.2em}

% 设置表格标题前后间距
\setlength{\abovecaptionskip}{0pt}
\setlength{\belowcaptionskip}{0pt}

% 代码格式自定义
\lstset{
    basicstyle=\sffamily,
    keywordstyle=\bfseries,
    commentstyle=\rmfamily\itshape,
    stringstyle=\ttfamily
}

% 命令重定义 
\renewcommand{\refname}{\bfseries{参~考~文~献}} %将Reference改为参考文献（用于 article）
% \renewcommand{\bibname}{参~考~文~献}%将bibiography改为参考文献（用于 book）
\renewcommand{\baselinestretch}{1.38}   %设置行间距
\renewcommand{\figurename}{\small\ttfamily 图}
\renewcommand{\tablename}{\small\ttfamily 表}


\title{我的xelatex文档模板}
\author{黎振胜}
\date{\today}

\begin{document}
\maketitle{}

\section{致谢和修改}
此文档以EvanYu的模板为基础，特此感谢！

修正了表格竖线不能闭合的问题，完善了sty文件和相关注释。

\section{说明}
之所以制作这个模板是因为用 \LaTeX{} 这么久了，每次要写一个新的文档都从头开始未免也太麻烦了，既然了解了那么多基础的知识了就做这么一个简单的 article 模板。虽然整理这个模板花了一些时间，不过这些时间都是相当的值得的，避免了我以后的重复性工作。这个文档是这个模板使用的一个示例文档，在 README.org 文件里面也含有这部分的源代码，便于写文档时参照使用。

\section{插入图片}
图片的插入对于文档来说是一个比较重要的方面，这里我选择的是插入一个浮动的环境，然后可以添加包括 caption、label 等标签，代码如下，示例见图\ref{pic:picturelabel}。

\begin{lstlisting}[language=tex]
\begin{figure}[h]
\centering
\label{pic:picturelabel}
\includegraphics[width=\textwidth]{./figure/9.eps}
\caption{\LaTeX{}图片插入}
\end{figure}
\end{lstlisting}

% 插入图片示例
\begin{figure}[h]
\centering
\label{pic:picturelabel}
\includegraphics[width=\textwidth]{./figure/9.eps}
\caption{\LaTeX{}图片插入}
\end{figure}

\section{插入表格}
\subsection{自动宽度表格}
示例的表格见表\ref{tab:table1}。自动宽度表格看似简单，但实际上包含的元素比较多。首先是使用 tabularx 来包含的，然后使用了多行和多列的单元格。其最重要的一点在于使用了自动列宽。第一列和第三列都是自动调整的宽度，而设置的表格宽度是 \textbackslash{}textwidth，除了第一列和第三列所用的宽度，其他的都留给了第二列（具体代码如下）。

\begin{verbatim}
\begin{table}[h]
\caption{指定宽度-列自动换行表格}\label{tab:table1}
\centering
\begin{tabularx}{\textwidth}{c|X|c}
\hline
\multirow{2}{*}{多行} & \multicolumn{2}{c}{多列} \\
\cline{2-3}
& data & comment \\
\hline
第一行 & 数据 & 注释 \\
\hline
\end{tabularx}
\end{table}
\end{verbatim}

% 指定表格宽度并设置自动换行示例
\begin{table}[h]
\caption{指定宽度-列自动换行表格}\label{tab:table1}
\centering
\begin{tabularx}{\textwidth}{c|X|c}
\hline
\multirow{2}{*}{多行} & \multicolumn{2}{c}{多列} \\
\cline{2-3}
& data & comment \\
\hline
第一行 & 数据 & 注释 \\
\hline
\end{tabularx}
\end{table}

\subsection{斜线表格}
斜线表格使用了 diagbox 宏包，使用这个宏包可以非常简单的做出斜线表，整个宏包只有一条命令 \textbackslash{}diagbox。两种使用方式，一种是跟两个参数，即划一条斜线；另一种是跟三个参数，即划两条斜线。具体的使用见如下的代码和表\ref{tab:table2}。

\begin{verbatim}
\begin{table}[h]
\caption{普通表格}\label{tab:table2}
\centering
\begin{tabular}{|l|ccc|}
\hline
\diagbox{名称}{编号} & id & data & comment \\
\hline
名称 & 1 & 数据 & 注释好像没有有\\
\hline
\end{tabular}
\end{table}
\end{verbatim}

% 普通表格居中并使用diagbox示例
\begin{table}[h]
\caption{普通表格}\label{tab:table2}
\centering
\begin{tabular}{|l|ccc|}
\hline
\diagbox{名称}{编号} & id & data & comment \\
\hline
名称 & 1 & 数据 & 注释好像没有\\
\hline
\end{tabular}
\end{table}

\subsection{带脚注的表格}
由于有的时候需要使用带脚注的表格，而对于 \LaTeX{} 来说并没有真正的脚注表格，所以使用 threeparttable 这个宏包可以制作很不错的脚注表格，表格包含在 threeparttalbe 环境内，然后在表格中需要脚注的地方使用 \textbackslash{}tnote\{number\}，然后在最后的 tablenotes 环境中使用 \textbackslash{}item [number]。表格的示例见表\ref{tab:table3}，代码见如下。

\begin{verbatim}
\begin{table}[h]
\caption{带脚注表格}\label{tab:table3}
\centering
\begin{threeparttable}
\begin{tabularx}{\textwidth}{c|X|c}
\hline
\multirow{2}{*}{多行} & \multicolumn{2}{|c|}{多列} \\
\cline{2-3}
 & data\tnote{1} & comment \\
\hline
第一行 & 数据 & 注释\tnote{2} \\
\hline
\end{tabularx}
\begin{tablenotes}
\item \footnotesize{[1] 第一个脚注是data}
\item \footnotesize{[2] 第二个脚注是注释(即\textbackslash{}tnote命令所在的位置)}
\end{tablenotes}
\end{threeparttable}
\end{table}
\end{verbatim}

% 真正的带脚注表格（并设置表格字体）
\begin{table}[h]
\caption{带脚注表格}\label{tab:table3}
\centering
\begin{threeparttable}
\begin{tabularx}{\textwidth}{c|X|c}
\hline
\multirow{2}{*}{多行} & \multicolumn{2}{c}{多列} \\
\cline{2-3}
 & data\tnote{1} & comment \\
\hline
第一行 & 数据 & 注释\tnote{2} \\
\hline
\end{tabularx}
\begin{tablenotes}
\item \footnotesize{[1] 第一个脚注是data}
\item \footnotesize{[2] 第二个脚注是注释(即\textbackslash{}tnote命令所在的位置)}
\end{tablenotes}
\end{threeparttable}
\end{table}

% 参考文献示例
\section{参考文献使用}
参考文献使用的样式为 plainnat，bib 文件所在的位置为 data/reference.bib。当然，参考文献可以使用多个文件。比如我的参考文献比较多，分类存在 1.bib、2.bib 这两个文件中
可以使用命令：\textbackslash{}bibliography\{./bib/1.bib,./bib/2.bib\}。
我在 data/reference.bib 文件中加入了一条记录，使用看我的右上角\cite{刘海洋2013latex}。
另外，我这里为了让参考文献在一页的开头使用了\textbackslash{}newpage 命令，空出来一页之后再添加参考文献。
为了让参考文献的标题居中，在插入参考文献之前添加了如下命令：

\verb|\titleformat{\section}{\centering\sanhao}{\chaptertitlename}{15pt}{}|

注意，xelatex生成参考文献页面比较麻烦，需要以下四个步骤，有时候还需要手动删除一些临时文件后重试

\begin{lstlisting}[language=bash]
    xelatex zstexart.tex
    bibtex zstexart.aux
    xelatex zstexart.tex
    xelatex zstexart.tex
\end{lstlisting}

\section{后记}
原本我在 github 上是建立了一个 repos 的，但是总是觉得其结构不是很好，本来打算新建一个 branch 来进行结构修改，但这本身就是一个比较小的工程，所以干脆重新来过算了。暂时来说优化了以前很多的结构，也修复了一些错误。便于以后再进行修改。

\newpage
\titleformat{\section}{\centering\sanhao}{\chaptertitlename}{15pt}{}
% % 指定参考文献样式和 bib 文件
\bibliographystyle{plainnat}
\bibliography{./bib/reference}
\end{document}
